% Judul penuh meluber 62,75 pt di kepala halaman; hanya mark header yang
% dipendekkan (sekali pakai, lihat \judulpendek di preamble-en.tex).
\judulpendek{Riemann Integrals}
\chapter{Riemann Integrals and the Fundamental Theorem of Calculus}

A program computes the area under $y=x^2$ on $[0,1]$ with 1000 slices of equal
width. If the height of each rectangle is taken at the left endpoint, the result is $0.3328335$; if at
the right endpoint, $0.3338335$. The difference is $0.001$, and there is no reason to prefer either one.
The correct value is $1/3$, which does not appear in either computation. Increasing the slices to a
million does not help: both answers still bracket the value, they never catch it.

What is needed is not a finer partition, but a way to read two approximations at
once. If the difference can be made smaller than any positive number,
then there is one number trapped between them, and that is the integral. This entire
chapter rests on that simple idea.

Chapter 12 taught us to take change apart into rates: from position to velocity,
from a function to its derivative. What is still missing is the way home. We know that the velocity at
second $t$ is $v(t)$, but we do not yet know how much distance was covered in one
minute. The intuition is clear: distance is the sum of velocity times small slices of time.
What is not clear is how to add up infinitely many such terms correctly.
This chapter answers that, and the answer is surprising: the way home is the antiderivative.

\begin{untukanak}
Imagine we want to know the area of a lake whose shape is irregular. There is no
formula for such a shape. All we have is transparent graph paper. We
lay that paper over the drawing of the lake, then count how many squares lie entirely
inside the water. The result is too small. We count again how many squares are
touched by water, even slightly. The result is too large. The true area lies
between those two numbers.

Now cut the squares into ten times smaller ones. The lower number rises, the upper number
falls, and the two grow closer. If the squares are very small, the two numbers are almost equal,
and we can state the area of the lake to any accuracy we like.

There is one condition that is easy to overlook. If the lake has very thin coral islands
scattered everywhere, the squares that touch them can keep the upper estimate far
from the lower estimate. Mathematicians call that condition integrability. In this chapter we
shall learn why two bracketing estimates can be trusted, and when they fail.
\end{untukanak}

\section{Partitions, Upper Sums, and Lower Sums}
\tingkat{basic}

We want to give meaning to the area under a graph without assuming that the area already
has a meaning: enclose the region with rectangles from inside and from outside, then
shrink the difference.

\begin{definisi}[Partitions and upper-lower sums]\index{partition}
Let $[a,b]$ be a closed interval. A \emph{partition} $\mathcal{P}$ of $[a,b]$
is a finite set of points
\[
a=x_0<x_1<\cdots<x_n=b .
\]
The width of the $i$-th slice is $\Delta x_i=x_i-x_{i-1}$, and the \emph{mesh} of that partition
is the widest slice width, $\|\mathcal{P}\|=\max_i\Delta x_i$.

Let $f$ be bounded on $[a,b]$. On the $i$-th slice write
\[
M_i=\sup\{f(x):x_{i-1}\le x\le x_i\},\qquad
m_i=\inf\{f(x):x_{i-1}\le x\le x_i\}.
\]
The \emph{upper sum} and the \emph{lower sum} of the partition $\mathcal{P}$ are
\begin{equation}
U(\mathcal{P},f)=\sum_{i=1}^{n}M_i\,\Delta x_i,
\qquad
L(\mathcal{P},f)=\sum_{i=1}^{n}m_i\,\Delta x_i .
\end{equation}
\end{definisi}

We use supremum and infimum, not the largest and smallest values, because $f$ need not
attain its extremes on a slice: the function equal to $x$ for $x<1$ and to $0$
at $x=1$ is bounded on $[0,1]$, but its largest value on the last slice does
not exist, whereas its supremum does, namely $1$.

The upper sum is always at least equal to the lower sum, because $M_i\ge m_i$ for
every $i$. For $f(x)=x^2$ on $[0,1]$ with four slices of equal width, the lower
sum is
\[
L=\frac{1}{4}\Bigl(0+\frac{1}{16}+\frac{4}{16}+\frac{9}{16}\Bigr)=\frac{7}{32},
\]
and the upper sum is $\frac{15}{32}$. Both bracket the value sought, with a difference of
$\frac14$ that is still large. If $f$ takes negative values, these sums compute signed
area, and the minus sign is part of the definition, not an error.

Figure~\ref{fig:rect-halus} shows two values of $n$ at once, with the right endpoint
as height. At $n=4$ the staircase is coarse and $U_4=0.46875$; at $n=16$ the serration is
fine and $U_{16}=0.3652343750$. The distance from the staircase to the line $1/3$ shrinks from $0.1354$
to $0.0319$.

\begin{figure}[htbp]
\centering
\begin{tikzpicture}
\begin{axis}[width=0.47\linewidth,height=4.6cm,
  axis lines=left,xmin=0,xmax=1.02,ymin=0,ymax=1.12,
  xtick={0,0.5,1},ytick={0,0.5,1},
  tick label style={font=\scriptsize},label style={font=\scriptsize},
  xlabel={$x$},
  title={\scriptsize $n=4$: $U_4=0.46875$}]
\pgfplotsinvokeforeach{1,2,3,4}{
  \draw[fill=layerteal!18,draw=layerteal!65,line width=0.4pt]
    (axis cs:{(#1-1)/4},0) rectangle (axis cs:{#1/4},{#1*#1/16});
}
\draw[densely dashed,layergrey,line width=0.6pt]
  (axis cs:0,0.33333) -- (axis cs:1,0.33333);
\addplot[thick,accent,domain=0:1,samples=80] {x^2};
\node[font=\scriptsize,text=layergrey,anchor=west] at (axis cs:0.06,0.40) {$1/3$};
\end{axis}
\end{tikzpicture}\hfill
\begin{tikzpicture}
\begin{axis}[width=0.47\linewidth,height=4.6cm,
  axis lines=left,xmin=0,xmax=1.02,ymin=0,ymax=1.12,
  xtick={0,0.5,1},ytick={0,0.5,1},
  tick label style={font=\scriptsize},label style={font=\scriptsize},
  xlabel={$x$},
  title={\scriptsize $n=16$: $U_{16}=0.3652343750$}]
\pgfplotsinvokeforeach{1,...,16}{
  \draw[fill=layerteal!18,draw=layerteal!65,line width=0.3pt]
    (axis cs:{(#1-1)/16},0) rectangle (axis cs:{#1/16},{#1*#1/256});
}
\draw[densely dashed,layergrey,line width=0.6pt]
  (axis cs:0,0.33333) -- (axis cs:1,0.33333);
\addplot[thick,accent,domain=0:1,samples=80] {x^2};
\node[font=\scriptsize,text=layergrey,anchor=west] at (axis cs:0.06,0.40) {$1/3$};
\end{axis}
\end{tikzpicture}
\caption{Rectangular approximation of $f(x)=x^{2}$ on $[0,1]$ with right-endpoint heights.
At $n=4$ the upper sum is $U_4=0.46875$, at $n=16$ the upper sum is $U_{16}=0.3652343750$,
and both are still above the true value $1/3=0.333\ldots$ (dashed line).
The difference $U_n-L_n$ is exactly $1/n$, so going from $n=4$ to $n=16$ cuts the distance from the
staircase to that line down to a quarter of it.}
\label{fig:rect-halus}
\end{figure}


There are two facts that will be used again and again. The first: refining a partition
does not make the estimate worse.

\begin{lemma}[Refinement]\index{refinement}
If $\mathcal{P}^{*}$ is a refinement of $\mathcal{P}$, meaning that
$\mathcal{P}\subseteq\mathcal{P}^{*}$, then
\[
L(\mathcal{P},f)\le L(\mathcal{P}^{*},f)\le U(\mathcal{P}^{*},f)\le U(\mathcal{P},f).
\]
\end{lemma}

\begin{proof}
It suffices to prove this for a refinement that adds one point $c$ inside the slice
$[x_{i-1},x_i]$, since a general refinement is obtained by repeating that step.
Let $m'$ and $m''$ be the infima of $f$ on $[x_{i-1},c]$ and on $[c,x_i]$. Both
intervals are contained in $[x_{i-1},x_i]$, so that $m'\ge m_i$ and $m''\ge m_i$.
Consequently
\[
m'(c-x_{i-1})+m''(x_i-c)\ \ge\ m_i(c-x_{i-1})+m_i(x_i-c)=m_i\,\Delta x_i ,
\]
so the lower sum does not decrease. The same argument with suprema gives $M'\le M_i$
and $M''\le M_i$, so that the upper sum does not increase.
\end{proof}

The second fact is stronger: two different partitions can still be compared, as long as what
is compared is a lower sum with an upper sum.

\begin{lemma}
For any two partitions $\mathcal{P}_1$ and $\mathcal{P}_2$ of $[a,b]$ we have
\[
L(\mathcal{P}_1,f)\le U(\mathcal{P}_2,f).
\]
\end{lemma}

\begin{proof}
Take $\mathcal{P}$ the common refinement, that is, the union of the points of the two partitions. By
the previous lemma, $L(\mathcal{P}_1,f)\le L(\mathcal{P},f)$ and
$U(\mathcal{P},f)\le U(\mathcal{P}_2,f)$. Since $L(\mathcal{P},f)\le U(\mathcal{P},f)$
always holds, that chain of inequalities gives the result asked for.
\end{proof}

That lemma ensures that every lower sum is a lower bound for every upper
sum. Hence the set of all lower sums is bounded above, and the set of all upper sums
is bounded below; the following two numbers always exist.

\begin{definisi}[Riemann integrable]\index{Riemann integrable}
For $f$ bounded on $[a,b]$, the \emph{lower integral} and the \emph{upper integral}
are defined as
\[
\underline{I}(f)=\sup_{\mathcal{P}}L(\mathcal{P},f),
\qquad
\overline{I}(f)=\inf_{\mathcal{P}}U(\mathcal{P},f),
\]
with the supremum and infimum taken over all partitions of $[a,b]$. Since
$\underline{I}(f)\le\overline{I}(f)$, the function $f$ is called \emph{Riemann integrable}
on $[a,b]$ when those two numbers are equal. Their common value is written
$\int_a^b f$ or $\int_a^b f(x)\,dx$.
\end{definisi}

Notice what is absent from that definition: there is no sample point and no
limit over $n$. The measure of success of an approximation is the difference between its two brackets,
not its closeness to a guess.

\kamus{partition}{partisi}{A way of cutting an interval into several slices, like
dividing a bar of chocolate with a knife, not measuring exactly how long each slice is.}

\kamus{upper sum}{jumlah atas}{An estimate that is deliberately inflated, like
weighing goods together with their box.}

\begin{bacakeras}
Two voices. Call the first \emph{Cutter}, the second \emph{Adder}.

Cutter: I have a strip of ribbon. I cut it into four.
Adder: I add up those four pieces. The result is too small.
Cutter: I cut again. So, eight.
Adder: The result goes up a little. Still too small.
Cutter: Sixteen.
Adder: Up again. Closer, not there yet.
Cutter: A hundred.
Adder: Very close.
Cutter: A thousand.
Adder: Closer still, still not there.

Cutter: When do we get there?
Adder: Never, if you only add slices. Stop chasing the number,
start chasing the bracket.
Cutter: What bracket?
Adder: Thin slices give a lower bound. Fat slices give an upper
bound. If those two bounds meet at one point, that point is the answer.
Cutter: So my slices are never enough?
Adder: Never. Slices only bring the two bounds close. What answers
is not the last slice, but the meeting of the two bounds.

Now repeat one line: cut thin, bracket tight, bounds meet.
\end{bacakeras}

\section{The Integrability Criterion}
\tingkat{intermediate}

The definition just given speaks of suprema and infima over all partitions, and there are infinitely
many partitions. The following criterion replaces them with a single inequality that
can be checked on one partition alone.

\begin{teorema}[Integrability criterion]\index{integrability criterion}
A bounded function $f$ is Riemann integrable on $[a,b]$ if and only if for
every $\varepsilon>0$ there is a partition $\mathcal{P}$ with
\[
U(\mathcal{P},f)-L(\mathcal{P},f)<\varepsilon .
\]
\end{teorema}

\begin{proof}
For every partition,
$L(\mathcal{P},f)\le\underline{I}(f)\le\overline{I}(f)\le U(\mathcal{P},f)$; the first inequality
comes from the definition of supremum, the last from the definition of infimum.

$(\Rightarrow)$ Suppose $f$ is integrable with value $I$, and take
$\varepsilon>0$. By the definition of supremum, there is $\mathcal{P}_1$ with
$L(\mathcal{P}_1,f)>I-\varepsilon/2$. By the definition of infimum, there is $\mathcal{P}_2$
with $U(\mathcal{P}_2,f)<I+\varepsilon/2$. Take $\mathcal{P}$ the common refinement
of both. The refinement lemma gives
\[
U(\mathcal{P},f)-L(\mathcal{P},f)
<\Bigl(I+\tfrac{\varepsilon}{2}\Bigr)-\Bigl(I-\tfrac{\varepsilon}{2}\Bigr)
=\varepsilon .
\]

$(\Leftarrow)$ Suppose that condition is met. For every $\varepsilon>0$ there is
$\mathcal{P}$ with $U(\mathcal{P},f)-L(\mathcal{P},f)<\varepsilon$, so that
\[
0\le\overline{I}(f)-\underline{I}(f)\le U(\mathcal{P},f)-L(\mathcal{P},f)<\varepsilon .
\]
The number $\overline{I}(f)-\underline{I}(f)$ is nonnegative and smaller than every
$\varepsilon>0$, so that number is zero. The two integrals are equal, and $f$ is integrable.
\end{proof}

The criterion turns a question about infinitely many partitions into a single
inequality. A direct consequence: if the difference is small on one partition, that difference
stays small on every refinement of it, since the refinement lemma can only
narrow it.

Before looking for classes of functions that pass the criterion, we record the properties of the integral
that will be used throughout. All of them are consequences of the definition of upper and lower sums, not new
assumptions.

\begin{proposisi}[Basic properties of the integral]\index{integral, basic properties}
Let $f$ and $g$ be Riemann integrable on $[a,b]$. Then:
(a) $f+g$ is integrable and $\int_a^b(f+g)=\int_a^b f+\int_a^b g$;
(b) if $f\le g$ on $[a,b]$ then $\int_a^b f\le\int_a^b g$;
(c) if $a<c<b$ then $f$ is integrable on $[a,c]$ and on $[c,b]$, and
$\int_a^b f=\int_a^c f+\int_c^b f$;
(d) $|f|$ is integrable and $\bigl|\int_a^b f\bigr|\le\int_a^b|f|$;
(e) for every constant $k$, $\int_a^b k\,dx=k(b-a)$.
\end{proposisi}

\begin{proof}
(a) On every slice, $\sup(f+g)\le\sup f+\sup g$ and $\inf(f+g)\ge\inf f+\inf g$.
Hence $U(\mathcal{P},f+g)\le U(\mathcal{P},f)+U(\mathcal{P},g)$ and
$L(\mathcal{P},f+g)\ge L(\mathcal{P},f)+L(\mathcal{P},g)$. Take $\varepsilon>0$ and
choose $\mathcal{P}$ that makes $U-L$ less than $\varepsilon/2$ for $f$ and $g$
at once. Then
\[
U(\mathcal{P},f+g)-L(\mathcal{P},f+g)
\le\bigl[U(\mathcal{P},f)-L(\mathcal{P},f)\bigr]
+\bigl[U(\mathcal{P},g)-L(\mathcal{P},g)\bigr]<\varepsilon ,
\]
so that $f+g$ is integrable. Its value is obtained by bracketing from both sides, since
$\varepsilon$ is arbitrary.

(b) If $f\le g$, then $M_i(f)\le M_i(g)$ on every slice, so that
$U(\mathcal{P},f)\le U(\mathcal{P},g)$ for all $\mathcal{P}$. Take the infimum of both
sides.

(c) If $\mathcal{P}$ contains $c$, the sum splits into the parts $[a,c]$ and
$[c,b]$, with $L(\mathcal{P},f)=L(\mathcal{P}_1,f)+L(\mathcal{P}_2,f)$ and the same
for $U$. The supremum gives
$\underline{I}(f)\le\underline{I}_{[a,c]}(f)+\underline{I}_{[c,b]}(f)$, the infimum
gives the opposite direction. For a partition without $c$, add $c$ first; the
refinement lemma preserves the inequality.

(d) Since $-|f|\le f\le|f|$, parts (b) and (a) give
$-\int|f|\le\int f\le\int|f|$. The integrability of $|f|$ follows from
$\bigl||f(x)|-|f(y)|\bigr|\le|f(x)-f(y)|$, which gives
$M_i(|f|)-m_i(|f|)\le M_i(f)-m_i(f)$.

(e) On every slice, $M_i=m_i=k$, so $U=L=k(b-a)$ for all partitions.
\end{proof}

Part (c) is the reason we may cut an interval at any point; without it the accumulation
function could not be defined cleanly.

\section{Classes of Integrable Functions}
\tingkat{intermediate}

The criterion in the previous section is a test. Now we look for classes of functions that
pass that test without having to find a special partition for each case. The two most
useful classes are continuous functions and monotone functions.

\begin{akibat}[Continuity guarantees integrability]
If $f$ is continuous on $[a,b]$, then $f$ is Riemann integrable on $[a,b]$.
\end{akibat}

\begin{proof}
Take $\varepsilon>0$. Since $[a,b]$ is compact, $f$ is uniformly continuous there (Chapter 11):
there is $\delta>0$ such that $|f(x)-f(t)|<\varepsilon/(b-a)$ when $|x-t|<\delta$. Choose
a partition with mesh less than $\delta$. On every slice, for all $x,t$ in
that slice we have $|f(x)-f(t)|<\varepsilon/(b-a)$, so that
$M_i-m_i\le\varepsilon/(b-a)$. Then
\[
U(\mathcal{P},f)-L(\mathcal{P},f)
=\sum_{i=1}^{n}(M_i-m_i)\,\Delta x_i
\le\frac{\varepsilon}{b-a}\sum_{i=1}^{n}\Delta x_i=\varepsilon .
\]
The integrability criterion closes the proof.
\end{proof}

Note the role of uniform continuity: one $\delta$ must work for all slices
at once, whereas ordinary continuity only gives a $\delta$ that depends on the
point. This is where Chapter 11 is used, not as decoration.

The second class does not require continuity: a monotone function may jump, as long as
its jumps are finite in number.

\begin{teorema}[Monotonicity guarantees integrability]
If $f$ is monotone on $[a,b]$, then $f$ is Riemann integrable on $[a,b]$.
\end{teorema}

\begin{proof}
Suppose $f$ is increasing. Divide $[a,b]$ into $n$ slices of equal width, $\Delta x=(b-a)/n$.
On the $i$-th slice, the smallest value is at the left endpoint and the largest value at the right
endpoint, so that $M_i-m_i=f(x_i)-f(x_{i-1})$. Then
\[
U(\mathcal{P},f)-L(\mathcal{P},f)
=\Delta x\sum_{i=1}^{n}\bigl(f(x_i)-f(x_{i-1})\bigr)
=\frac{(b-a)\bigl(f(b)-f(a)\bigr)}{n},
\]
because the sum in the middle telescopes. Choose $n>(b-a)(f(b)-f(a))/\varepsilon$ and the criterion
is satisfied. If $f$ is decreasing, apply this result to $-f$ and use linearity.
\end{proof}

That proof uses no $\delta$: all it uses is that the sum of all the jumps does not
exceed $f(b)-f(a)$, so a monotone function is easy to integrate even though its graph is full of
corners.

Those two classes can be combined: a function that is continuous except at finitely many points
is still integrable (Rudin, Theorem 6.10, not repeated here).

Now the other side: bounded functions that fail. The cleanest example is the
Dirichlet function, equal to $1$ at rational numbers and to $0$ at irrational numbers.

\begin{contoh}[The Dirichlet function is not integrable]\index{Dirichlet function}
Let $D(x)=1$ if $x$ is rational and $D(x)=0$ if $x$ is irrational, on $[0,1]$.
Every slice contains both (Chapter 3, density of $\mathbb{Q}$), so that $M_i=1$ and
$m_i=0$ for every $i$, and
\[
U(\mathcal{P},D)=\sum_{i=1}^{n}1\cdot\Delta x_i=1,
\qquad
L(\mathcal{P},D)=0
\]
for \emph{every} partition. The difference is always $1$, never smaller than
$\varepsilon=1/2$, so by the criterion $D$ is not integrable even though it is bounded.
\end{contoh}

That failure is not because $D$ is broken at many points, but because
its jumps are spread everywhere. The following function shows that the number of
discontinuity points is not an obstacle.

\begin{teorema}[Thomae's function]\index{Thomae function}
Define $T$ on $[0,1]$ by $T(x)=1/q$ if $x=p/q$ in lowest
terms with $q\ge1$, and $T(x)=0$ if $x$ is irrational. Then $T$ is integrable in the
Riemann sense and $\int_0^1 T=0$.
\end{teorema}

\begin{proof}
Take $\varepsilon>0$. The set $S=\{x:T(x)\ge\varepsilon\}$ consists of fractions with
denominator $q\le1/\varepsilon$; for each $q$ there are only $q+1$ such fractions in
$[0,1]$, so that $S$ is finite. Cover $S$ with finitely many open intervals
whose total length is less than $\varepsilon$.

Form a partition that contains all endpoints of those covering intervals. The slices that contain $S$
have total length less than $\varepsilon$ and there $M_i-m_i\le1$, so
their contribution is less than $\varepsilon$. The other slices satisfy $M_i\le\varepsilon$ and
$m_i\ge0$, so their contribution is at most $\varepsilon$. Hence
$U(\mathcal{P},T)-L(\mathcal{P},T)<2\varepsilon$, and the criterion closes the proof. The value is
zero because every slice contains an irrational number, so that $L(\mathcal{P},T)=0$.
\end{proof}

So $T$ is broken at every rational number, which is dense in $[0,1]$, yet $T$ is still
integrable. What decides is not the number of discontinuity points but the amount of space
they occupy: Lebesgue's criterion states that a bounded function is Riemann integrable
if and only if its discontinuity points have measure zero (Rudin, Theorem 11.33(b); Abbott,
Chapter 7). We quote it without proof because measure demands measure theory.

\begin{tantangan}
Define $h$ on $[0,1]$ by $h(x)=1$ if $x=p/q$ in lowest terms
and the denominator $q$ is odd, and $h(x)=0$ for all other $x$, including every irrational
number. Is $h$ Riemann integrable on $[0,1]$? If so, what is the value of
the integral? If not, prove it using the integrability criterion, not by
counting points.

Hint 1: do not count the number of points where $h$ equals one. The number is
infinite, and the criterion does not ask for that count.

Hint 2: take any slice $[x_{i-1},x_i]$ with $x_{i-1}<x_i$. That slice
contains fractions with odd denominator, so $M_i=1$, and it contains irrational numbers, so
$m_i=0$. Thus $U-L=1$ for every partition.

If you get stuck, go back to the section Classes of Integrable Functions, the Dirichlet example.
\end{tantangan}

\kamus{Riemann integrability}{keterintegralan}{The property of a function whose lower and
upper estimates can be brought together until their difference vanishes.}

\begin{untukanak}
Have you ever looked at a car's speedometer and then asked how much distance has been
covered? Distance is not displayed there, only instantaneous speed. But if we
record the speed every second, we can add it up: one second at
ten meters per second means 10 meters, the next second 12 meters, and
so on. The sum of all those slices is the distance.

The problem is that the speed changes in the middle of a second. A car can speed up and then slow down
within the same second. If we use the speed at the start of the second, the distance we
compute is too small. If we use the speed at the end of the second, the result is too
large. The right value lies between them.

If we record ten times per second, the two estimates grow closer; a hundred times,
closer still. Mathematicians ask further: if the recording can be as dense as we
like, do the two estimates head toward one and the same number? The answer is not always yes,
and in this chapter we shall learn when it is yes and why.
\end{untukanak}

\section{The Integral as an Accumulation Function}
\tingkat{intermediate}

Up to here the integral is only a number attached to an interval. We can
view it as a function by letting the upper endpoint move: the integral becomes
a report on how much has accumulated up to that point.

\begin{definisi}[Accumulation function]\index{accumulation function}
Let $f$ be integrable on $[a,b]$. For $x\in[a,b]$ define
\[
F(x)=\int_a^x f(t)\,dt .
\]
For $x<a$ we use the convention $\int_a^x f=-\int_x^a f$, and for $x=a$ the value is
zero. The function $F$ is called the \emph{accumulation function} of $f$ on $[a,b]$.
\end{definisi}

With that convention the additive property holds for any ordering of points, so that
\begin{equation}
F(y)-F(x)=\int_x^{y} f(t)\,dt
\end{equation}
for all $x,y\in[a,b]$. This is the form used in all the proofs that follow:
the increase of savings between two times is the integral of the income rate in between.

The accumulation function is always better than the function it comes from: it is continuous even when $f$
is not, and differentiable at every point where $f$ is continuous.

\begin{proposisi}[Accumulation is always continuous]
If $f$ is bounded on $[a,b]$ with $|f|\le M$ and integrable, then the $F$ that
is defined above satisfies
\[
|F(y)-F(x)|\le M\,|y-x|
\]
for all $x,y\in[a,b]$. In particular $F$ is continuous on $[a,b]$.
\end{proposisi}

\begin{proof}
For $x<y$, monotonicity and part (d) of the basic properties proposition give
\[
|F(y)-F(x)|=\Bigl|\int_x^{y}f\Bigr|\le\int_x^{y}|f|\le M(y-x).
\]
For $y<x$ the same result with the opposite sign. Hence $F$ satisfies a
Lipschitz condition with constant $M$, and every Lipschitz function is continuous.
\end{proof}

Now the first direction of the derivative-integral relationship. In everyday words: the rate of
growth of savings at a given moment equals the income rate at that moment.

\begin{teorema}[Fundamental theorem of calculus, first part]\index{fundamental theorem of calculus}
Let $f$ be integrable on $[a,b]$ and continuous at the point $x_0\in(a,b)$. Then the
accumulation function $F$ is differentiable at $x_0$ and
\[
F'(x_0)=f(x_0).
\]
\end{teorema}

\begin{proof}
Take $h\ne0$ with $x_0+h\in[a,b]$. From the increment form,
\[
\frac{F(x_0+h)-F(x_0)}{h}-f(x_0)
=\frac{1}{h}\int_{x_0}^{x_0+h}\bigl(f(t)-f(x_0)\bigr)\,dt ,
\]
because the integral of the constant function $f(x_0)$ over an interval of length $|h|$ equals
$f(x_0)h$.

Take $\varepsilon>0$. Since $f$ is continuous at $x_0$, there is $\delta>0$ such that
$|t-x_0|<\delta$ implies $|f(t)-f(x_0)|<\varepsilon$. For $0<|h|<\delta$ the whole
interval of integration lies inside $(x_0-\delta,x_0+\delta)$, so that, by monotonicity
and part (d) of the basic properties proposition,
\[
\left|\frac{1}{h}\int_{x_0}^{x_0+h}\bigl(f(t)-f(x_0)\bigr)dt\right|
\le\frac{1}{|h|}\left|\int_{x_0}^{x_0+h}\bigl|f(t)-f(x_0)\bigr|\,dt\right|
\le\frac{1}{|h|}\,\varepsilon\,|h|=\varepsilon .
\]
Since $\varepsilon$ is arbitrary, the difference quotient tends to $f(x_0)$.
\end{proof}

That theorem demands continuity at only one point. Where $f$ jumps, the derivative of
$F$ may fail to exist, and that is not a weakness of the proof: for $f(t)=\operatorname{sgn} t$
on $[-1,1]$ we have $F(x)=\int_0^x\operatorname{sgn} t\,dt=|x|$, which is not
differentiable at $0$, exactly where the sign of $f$ changes.

One important consequence: every continuous function has an antiderivative, and its antiderivative can be
written as an integral. Chapter 12 derived the differentiation rules but did not answer whether
every continuous function is the derivative of something. Now that question is answered,
and the answer is constructive.

\begin{akibat}[An antiderivative always exists]
If $f$ is continuous on $[a,b]$, then $F(x)=\int_a^x f$ is differentiable on $[a,b]$
with $F'=f$. Hence $F$ is an antiderivative of $f$ on $[a,b]$.
\end{akibat}

\begin{proof}
The continuity of $f$ gives integrability (the previous corollary) and continuity at every
point, so that the first part of the theorem applies at every $x_0\in(a,b)$. At the endpoints
one-sided derivatives are used.
\end{proof}

The chain rule gives one further flexibility: if the upper limit depends on $x$,
\[
\frac{d}{dx}\int_a^{g(x)}f(t)\,dt=f\bigl(g(x)\bigr)\,g'(x),
\]
provided $g$ is differentiable. For example,
$\frac{d}{dx}\int_0^{x^2}e^{-t^2}dt=2x\,e^{-x^4}$.

The two directions are easy to confuse, so it is worth drawing them once.
Figure~\ref{fig:tdk-dua-arah} places both in one diagram: the upper arrow
accumulates thin slices, the lower arrow differentiates. What the theorem guarantees is that
differentiating the result of accumulating returns the original function, not a tidy algebraic inverse,
since an antiderivative is fixed only up to a constant.

\begin{figure}[htbp]
\centering
\begin{tikzpicture}[>=Stealth,every node/.style={font=\footnotesize},
  kotak/.style={draw,rounded corners=3pt,align=center,inner sep=5pt,
    minimum height=10mm,text width=3.3cm}]
\node[kotak,fill=layerindigo!8,draw=layerindigo] (f) at (0,0) {function $f$};
\node[kotak,fill=layerbrown!8,draw=layerbrown] (F) at (7.6,0)
  {accumulation function $F$};
\draw[->,line width=0.9pt,draw=layerindigo] ([yshift=2mm]f.north) to[bend left=24]
  node[above,text=layerindigo,align=center]
  {definite integral: accumulate thin slices} ([yshift=2mm]F.north);
\draw[->,line width=0.9pt,draw=layerbrown] ([yshift=-2mm]F.south) to[bend left=24]
  node[below,text=layerbrown,align=center]
  {derivative: $F'=f$, this is part 1} ([yshift=-2mm]f.south);
\node[draw=layergrey,dashed,rounded corners=3pt,align=center,inner sep=4pt]
  at (3.8,-2.7)
  {part 2 closes the loop: $\int_a^b F'(x)\,dx=F(b)-F(a)$};
\end{tikzpicture}
\caption{The two directions of the fundamental theorem of calculus as back-and-forth arrows. Upper arrow:
accumulating thin slices produces $F(x)=\int_a^x f$. Lower arrow: differentiating $F$
returns $f$, and that is the first part. The dashed box below is the second
part, which uses an antiderivative to compute integrals. For $f(x)=x^2$ on
$[0,1]$, $F(x)=x^3/3$, and the second part gives $F(1)-F(0)=1/3$, the same number as
the bracketing in Figure~\ref{fig:rect-halus}.}
\label{fig:tdk-dua-arah}
\end{figure}

\section{The Fundamental Theorem of Calculus, Second Part}
\tingkat{advanced}

The first part runs from the integral to the derivative. The second part runs the other way:
if we already hold an antiderivative, the integral can be computed with one subtraction.
In everyday words: total change equals the final value minus the initial value.

\begin{teorema}[Fundamental theorem of calculus, second part]
Let $F$ be continuous on $[a,b]$, differentiable on $(a,b)$, and let $F'$ be
Riemann integrable on $[a,b]$. Then
\[
\int_a^b F'(x)\,dx=F(b)-F(a).
\]
\end{teorema}

\begin{proof}
Take a partition $\mathcal{P}=\{x_0,\dots,x_n\}$ of $[a,b]$. On every subinterval, $F$ is
continuous on $[x_{i-1},x_i]$ and differentiable inside it, so the mean value
theorem (Chapter 12) gives $t_i\in(x_{i-1},x_i)$ with
\[
F(x_i)-F(x_{i-1})=F'(t_i)\,\Delta x_i .
\]
Since $m_i\le F'(t_i)\le M_i$, summing over $i$ gives
\[
L(\mathcal{P},F')\le\sum_{i=1}^{n}\bigl(F(x_i)-F(x_{i-1})\bigr)\le U(\mathcal{P},F').
\]
The sum in the middle telescopes and equals $F(b)-F(a)$, so that number is trapped
between the lower sum and the upper sum of every partition. Since $F'$ is integrable, both
brackets head toward the same value, namely $\int_a^b F'$.
\end{proof}

\begin{resep}
\begin{enumerate}
\item Find $F$ with $F'=f$, then check by differentiating $F$. If the derivative does not
      return to $f$, the candidate is wrong and the following steps are wasted.
\item Write $F$ at the upper limit, write $F$ at the lower limit, then subtract:
      $F(b)-F(a)$. Do not insert $+C$; the constant cancels.
\item Fix the sign if the interval is reversed: $\int_b^a f=-\int_a^b f$, and
      $\int_a^a f=0$.
\item Test with one coarse slice, a trapezoid or \texttt{quad}. If the result is far off,
      the error is in step 1 or 2, not in step 3.
\end{enumerate}
Warning about order: checking the antiderivative must happen before subtracting. Two values of an
arbitrary function still give a tidy and plausible difference, so that a mistake in
step 1 is not visible from the shape of the answer.
\end{resep}

Likewise the two parts do not replace each other: the first part differentiates the
accumulation function without requiring an antiderivative, whereas the second part requires an
antiderivative and integrability of its derivative.

The second part explains why the result is called net change, not area. For
$F(x)=\frac{x^4}{4}-\frac{x^2}{2}$ on $[-1,1]$, $F'(x)=x^3-x$, so that
\[
\int_{-1}^{1}(x^3-x)\,dx=F(1)-F(-1)=\Bigl(\frac14-\frac12\Bigr)-\Bigl(\frac14-\frac12\Bigr)=0 .
\]
The region above the axis has area $1/4$, the region below the axis also $1/4$, so the total area is
$1/2$ whereas the integral is zero because the signs are opposite.
Figure~\ref{fig:luas-bertanda} shows both.

\begin{figure}[htbp]
\centering
\begin{tikzpicture}
\begin{axis}[width=0.74\linewidth,height=5.4cm,
  axis lines=middle,xmin=-1.3,xmax=1.3,ymin=-0.6,ymax=0.6,
  xtick={-1,0,1},ytick=\empty,
  tick label style={font=\scriptsize},
  xlabel={$x$},ylabel={$y=x^3-x$},
  xlabel style={font=\scriptsize},ylabel style={font=\scriptsize,rotate=-90}]
\addplot[fill=layerjade!22,draw=layerjade!70,domain=-1:0,samples=60]
  {x^3-x} \closedcycle;
\addplot[fill=layerplum!22,draw=layerplum!70,domain=0:1,samples=60]
  {x^3-x} \closedcycle;
\addplot[thick,ink,domain=-1.2:1.2,samples=120] {x^3-x};
\node[font=\scriptsize,text=layerjade!55!black] at (axis cs:-0.52,0.36)
  {$+\frac14$};
\node[font=\scriptsize,text=layerplum!55!black] at (axis cs:0.52,-0.36)
  {$-\frac14$};
\end{axis}
\end{tikzpicture}
\caption{Signed area for $F'(x)=x^3-x$ on $[-1,1]$. The region above the $x$-axis
contributes $+1/4$, the region below the axis contributes $-1/4$, so that
$\int_{-1}^{1}(x^3-x)\,dx=F(1)-F(-1)=0$. The geometric area is still $1/2$; what is zero is
the signed sum, and that is what the second part of the theorem computes.}
\label{fig:luas-bertanda}
\end{figure}

The integrability hypothesis in the second part is not excessive. There are
differentiable functions whose derivative is bounded but not Riemann integrable; the classical
example was found by Volterra in the nineteenth century, and we do not
prove it here. The second part of the theorem is not an unconditional statement.

\section{Substitution and Integration by Parts}
\tingkat{advanced}

The two most frequently used rules for computing integrals are not new axioms: both are
differentiation rules read backwards, and their proofs use only the fundamental theorem of calculus.
The first one reverses the chain rule.

\begin{proposisi}[Substitution]\index{substitution}
Let $\varphi$ be differentiable with $\varphi'$ continuous on $[a,b]$, and let $f$ be
continuous on an interval that contains $\varphi([a,b])$. Then
\[
\int_a^b f\bigl(\varphi(t)\bigr)\,\varphi'(t)\,dt
=\int_{\varphi(a)}^{\varphi(b)}f(u)\,du .
\]
\end{proposisi}

\begin{proof}
Since $f$ is continuous, it has an antiderivative $G(u)=\int_c^u f$, with $c$ inside the
range of $\varphi$; by the first part of the theorem, $G'=f$. The chain rule (Chapter 12)
gives
\[
\frac{d}{dt}G\bigl(\varphi(t)\bigr)=G'\bigl(\varphi(t)\bigr)\varphi'(t)
=f\bigl(\varphi(t)\bigr)\varphi'(t).
\]
The right-hand side is continuous, hence integrable. The second part of the theorem applied to the function
$G\circ\varphi$ gives
\[
\int_a^b f\bigl(\varphi(t)\bigr)\varphi'(t)\,dt
=G\bigl(\varphi(b)\bigr)-G\bigl(\varphi(a)\bigr)
=\int_{\varphi(a)}^{\varphi(b)}f(u)\,du .
\]
\end{proof}

The second rule reverses the product rule.

\begin{proposisi}[Integration by parts]\index{integration by parts}
Let $F$ and $G$ be differentiable on $[a,b]$ with $F'$ and $G'$ continuous. Then
\[
\int_a^b F(x)\,G'(x)\,dx=F(b)G(b)-F(a)G(a)-\int_a^b F'(x)\,G(x)\,dx .
\]
\end{proposisi}

\begin{proof}
Let $H=FG$. The product rule gives $H'=F'G+FG'$, and the right-hand side is continuous
so that $H'$ is integrable. The second part of the theorem applied to $H$ gives
\[
F(b)G(b)-F(a)G(a)=\int_a^b H'=\int_a^b F'G+\int_a^b FG' .
\]
Move the first term to the left-hand side.
\end{proof}

\kamus{antiderivative}{antiturunan}{A function whose derivative is the function we are
integrating, like the way home on a road: it is useful only when the way out is already
known.}

As a numerical check, take $\int_0^1 x\,e^{x^2}dx$. The substitution $u=x^2$,
$du=2x\,dx$ gives
\[
\int_0^1 x\,e^{x^2}dx=\frac12\int_0^1 e^{u}du=\frac{e-1}{2}=0.8591409142295225 .
\]
\texttt{scipy.integrate.quad} gives $0.8591409142295227$, a difference of $2\times10^{-16}$
from machine rounding. For the second rule, $F(x)=x$ and $G'(x)=e^x$ give
$\int_0^1 xe^xdx=e-(e-1)=1$, and \texttt{quad} returns $1.0$ with estimated
error $1.1\times10^{-14}$.

\section{The Mean Value Theorem for Integrals}
\tingkat{intermediate}

An average is usually computed by adding and then dividing by the number of data. For a function
that changes continuously, the sum is replaced by an integral and the divisor by the length of the interval: the average
is the height of the rectangle whose area equals the area under the curve.

\begin{definisi}[Average value]
For $f$ integrable on $[a,b]$, the \emph{average value} of $f$ on that interval
is
\[
\bar{f}=\frac{1}{b-a}\int_a^b f(x)\,dx .
\]
\end{definisi}

\begin{teorema}[Mean value for integrals]\index{mean value of an integral}
If $f$ is continuous on $[a,b]$, then there is $c\in[a,b]$ with
\[
f(c)=\frac{1}{b-a}\int_a^b f(x)\,dx .
\]
\end{teorema}

\begin{proof}
Since $f$ is continuous on a closed bounded interval, it attains its smallest value $m$ and
its largest value $M$ there (Chapter 11). Monotonicity gives
\[
m(b-a)\le\int_a^b f\le M(b-a),
\]
so that $\bar{f}\in[m,M]$. By the intermediate value theorem (Chapter 11), the continuous function $f$
takes every value between $m$ and $M$, so there is $c\in[a,b]$ with
$f(c)=\bar{f}$.
\end{proof}

As a number, for $f(x)=x^2$ on $[0,1]$ the average is $1/3$, and the point that the
theorem reaches is $c=1/\sqrt3=0.5773502691896258$, since
$c^2=0.3333333333333334$. That point is not the midpoint of the interval: the theorem guarantees
existence, not location.

\section{Worked Examples}
\tingkat{intermediate}

Now we work through three things: the upper and lower sums of $f(x)=x^2$ on $[0,1]$,
a comparison of five quadrature rules, and one integral without an elementary antiderivative.
Fractions are computed with \texttt{fractions.Fraction}, without rounding.

Start from $n$ slices of equal width with $x_i=i/n$ and $\Delta x=1/n$, so that
\[
U_n=\sum_{i=1}^{n}\Bigl(\frac{i}{n}\Bigr)^2\frac1n=\frac{1}{n^3}\sum_{i=1}^{n}i^2
=\frac{(n+1)(2n+1)}{6n^2}=\frac13+\frac{1}{2n}+\frac{1}{6n^2},
\]
using the sum-of-squares formula from Chapter 1. In the same way
$L_n=\frac13-\frac{1}{2n}+\frac{1}{6n^2}$, so that
\begin{equation}
U_n-L_n=\frac1n .
\end{equation}
That result is not an approximation. For every $n$ the difference is exactly $1/n$; since $1/n$ can
be made smaller than any $\varepsilon$ by taking $n>1/\varepsilon$, the integrability
criterion is satisfied. Hence the following table is at the same time a proof that
$\int_0^1x^2dx=1/3$: the two sums bracket $1/3$, and their distance tends to zero.

\begin{center}\small
\begin{tabular}{rrrr}
\hline
$n$ & $U_n$ & $L_n$ & $U_n-L_n$ \\
\hline
1 & $1.0000000000$ & $0.0000000000$ & $1$ \\
4 & $0.4687500000$ & $0.2187500000$ & $1/4$ \\
16 & $0.3652343750$ & $0.3027343750$ & $1/16$ \\
64 & $0.3411865234$ & $0.3255615234$ & $1/64$ \\
256 & $0.3352890015$ & $0.3313827515$ & $1/256$ \\
1024 & $0.3338217735$ & $0.3328452110$ & $1/1024$ \\
\hline
\end{tabular}
\end{center}

At $n=1024$ the difference is $9.77\times10^{-4}$ and the distance from $U_n$ to $1/3$ is half of that.
Each time $n$ is doubled the difference is exactly halved, and that is why upper and lower
sums are not an efficient way to reach high accuracy.

\begin{tangan}
Materials: one sheet of graph paper, scissors, a ruler, tape, a pencil.

Steps. (1) On the graph paper draw a curve that rises faster and faster: one square from the
left edge has height one square, two squares have height four squares, three squares have height
nine squares. (2) Cut the region under the curve into vertical strips one square wide.
(3) Lay those strips side by side on the table, from shortest to longest, with level
bases. (4) Look at their upper edge: what appears is a jagged staircase, not a smooth line.
(5) Repeat with strips half a square wide, then a quarter square wide.

What we should see: the serration on the upper edge grows finer as the strips grow thinner,
whereas the total area does not change, since cutting and rearranging neither adds nor removes
paper. Hence the area of that curved region equals the area of the jagged figure, and that jagged
figure increasingly resembles one rectangle. The height of that rectangle is what we shall later call
the average value of the function on that interval.
\end{tangan}

Now compare five quadrature rules on the same integral: left endpoint, right
endpoint, midpoint $x_i=(i-\frac12)/n$, trapezoid, and Simpson with weights
$1,4,2,4,\dots,4,1$ divided by $3n$. The error is computed against the exact value $1/3$.

\begin{center}\small
\begin{tabular}{rrrrrr}
\hline
$n$ & left & right & midpoint & trapezoid & Simpson \\
\hline
10 & $4.833\times10^{-2}$ & $5.167\times10^{-2}$ & $8.333\times10^{-4}$
   & $1.667\times10^{-3}$ & $1.1\times10^{-16}$ \\
100 & $4.983\times10^{-3}$ & $5.017\times10^{-3}$ & $8.333\times10^{-6}$
   & $1.667\times10^{-5}$ & $0$ \\
1000 & $4.998\times10^{-4}$ & $5.002\times10^{-4}$ & $8.333\times10^{-8}$
   & $1.667\times10^{-7}$ & $0$ \\
\hline
\end{tabular}
\end{center}

Three patterns stand out. The error of the left and right rules drops by a factor of ten each time $n$ grows
by a factor of ten: order $\mathcal{O}(1/n)$, about $1/(2n)$ because $f$ is increasing. The error of the midpoint
and trapezoid rules drops by a factor of a hundred, order $\mathcal{O}(1/n^2)$, with values
$1/(12n^2)$ and $1/(6n^2)$; the two have opposite signs and the trapezoid is twice as large,
so that $(2M+T)/3$ cancels the term. Simpson is exact to machine rounding, since the
rule is exact for polynomials of degree at most three.

The third integral has no elementary antiderivative:
\[
\int_0^1 e^{-x^2}dx .
\]
That absence is a consequence of Liouville's theorem, which we quote without proof. The reference
is the error function $\frac{\sqrt{\pi}}{2}\operatorname{erf}(1)=0.746824132812427$; \texttt{quad} gives
the same number with estimated error $8.3\times10^{-15}$. Simpson's rule gives an error of
$3.563\times10^{-4}$ at $n=2$, $1.988\times10^{-6}$ at $n=8$, and
$7.795\times10^{-9}$ at $n=32$.

\kertas{For $f(x)=x^2$ on $[0,1]$, derive the formula
$U_n-L_n=1/n$ yourself from the sum-of-squares formula, then compute $U_{10}$ and $L_{10}$ as
fractions. Check that their average $(U_{10}+L_{10})/2$ equals
$1/3+1/(6\cdot100)$, not $1/3$. Explain why that average is still not the best
estimate, even though its error is far smaller than the error of $U_{10}$ itself.}

\section{Improper Integrals and Their Relation to Series}
\tingkat{advanced}

The interval of integration has so far always been bounded and the function bounded. Many real
computations are not like that: we want the area of a region that extends without end, or the integral
of a function that blows up at the endpoint of the interval. The method is not to widen the Riemann
definition, but to write that integral as a limit of ordinary integrals.

\begin{definisi}[Improper integral]\index{improper integral}
If $f$ is Riemann integrable on $[a,R]$ for every $R>a$, define
\[
\int_a^\infty f=\lim_{R\to\infty}\int_a^R f ,
\]
provided that limit exists and is finite. If $f$ is integrable on $[c,b]$ for every
$c\in(a,b]$ but is unbounded near $a$, define
\[
\int_a^b f=\lim_{c\to a^{+}}\int_c^b f .
\]
In both cases, the integral is called \emph{convergent} if the limit exists, and
\emph{divergent} if not.
\end{definisi}

The most important example is the family $f(x)=x^{-p}$ on $[1,\infty)$. For $p\ne1$,
\[
\int_1^R x^{-p}dx=\frac{R^{1-p}-1}{1-p},
\]
so that the limit exists only when $p>1$, with value $1/(p-1)$; for $p=1$ the result is
$\ln R$, which tends to infinity. Hence $\int_1^\infty x^{-2}dx=1$, and the partial
integrals head toward that value at rate $1/R$: $0.9$, $0.99$, $0.999$ for
$R=10,100,1000$.

The link with the series of Chapter 5 is no coincidence: for $f$ positive and decreasing, the areas of
rectangles of height $f(n)$ bracket the area under the curve, so that
\[
\sum_{n=N+1}^{\infty}f(n)\ \le\ \int_N^\infty f(x)\,dx\ \le\ \sum_{n=N}^{\infty}f(n) .
\]
Hence the integral test holds: $\sum f(n)$ converges if and only if
$\int_1^\infty f$ converges. For $f(x)=1/x^2$ both converge, with
$\sum_{n=1}^{\infty}1/n^2=\pi^2/6=1.6449340668482264$. The left inequality also gives
an error bound, and the larger $N$ is the closer the error comes to that bound:

\begin{center}\small
\begin{tabular}{rrrr}
\hline
$N$ & error of the partial sum & $1/N$ \\
\hline
10 & $9.516634\times10^{-2}$ & $0.1$ \\
100 & $9.950167\times10^{-3}$ & $0.01$ \\
1000 & $9.995002\times10^{-4}$ & $0.001$ \\
\hline
\end{tabular}
\end{center}

It must be stressed: an improper integral is not a Riemann integral, but its limit. Hence
substitution and integration by parts apply only when the limit exists, and for functions
that change sign their convergence need not be absolute.

\section{Computational Lab}
\tingkat{advanced}

The following code is run with Python 3, \texttt{scipy} 1.18.0, and the module
\texttt{fractions} for exact arithmetic. All partitions are of equal width on $[0,1]$.

\begin{lab}
\begin{verbatim}
import math
from fractions import Fraction as F
from scipy.integrate import quad

def both_sums(n):
    U = sum(F(i * i, n ** 3) for i in range(1, n + 1))
    L = sum(F(i * i, n ** 3) for i in range(n))
    return U, L

print("n        U_n          L_n          U_n - L_n")
for n in (1, 2, 4, 8, 16, 32, 64, 128, 256, 512, 1024):
    U, L = both_sums(n)
    print(f"{n:5d}  {float(U):.10f}  {float(L):.10f}  {float(U - L):.2e}")

def rules(f, n):
    h = 1 / n
    left = sum(f(i * h) for i in range(n)) * h
    right = sum(f(i * h) for i in range(1, n + 1)) * h
    mid = sum(f((i + 0.5) * h) for i in range(n)) * h
    trap = (f(0) / 2 + sum(f(i * h) for i in range(1, n)) + f(1) / 2) * h
    simp = (f(0) + f(1) + sum((4 if i % 2 else 2) * f(i * h)
                              for i in range(1, n))) * h / 3
    return left, right, mid, trap, simp

f = lambda x: x * x
print()
print("error |approximation - 1/3|, integral of x^2 on [0,1]")
print("   n      left     right       mid  trapezoid    Simpson")
for n in (10, 100, 1000):
    err = [f"{abs(v - 1 / 3):.3e}" for v in rules(f, n)]
    print(f"{n:4d}  " + "  ".join(err))

g = lambda x: math.exp(-x * x)
v, e = quad(g, 0, 1)
ref = math.sqrt(math.pi) / 2 * math.erf(1)
print()
print(f"quad exp(-x^2) = {v:.15f}   estimated error {e:.1e}")
print(f"erf reference  = {ref:.15f}   difference {v - ref:+.1e}")
print(f"quad x^2       = {quad(f, 0, 1)[0]:.17f}")
print()
print("error |Simpson - reference| for exp(-x^2)")
for n in (2, 4, 8, 16, 32):
    print(f"  n={n:3d}  {abs(rules(g, n)[4] - ref):.3e}")
print()
print("error |trapezoid - reference| for exp(-x^2)")
for n in (10, 100, 1000):
    print(f"  n={n:4d}  {abs(rules(g, n)[3] - ref):.3e}")
\end{verbatim}
\end{lab}

The output obtained, as it stands:

\begin{verbatim}
n        U_n          L_n          U_n - L_n
    1  1.0000000000  0.0000000000  1.00e+00
    2  0.6250000000  0.1250000000  5.00e-01
    4  0.4687500000  0.2187500000  2.50e-01
    8  0.3984375000  0.2734375000  1.25e-01
   16  0.3652343750  0.3027343750  6.25e-02
   32  0.3491210938  0.3178710938  3.12e-02
   64  0.3411865234  0.3255615234  1.56e-02
  128  0.3372497559  0.3294372559  7.81e-03
  256  0.3352890015  0.3313827515  3.91e-03
  512  0.3343105316  0.3323574066  1.95e-03
 1024  0.3338217735  0.3328452110  9.77e-04

error |approximation - 1/3|, integral of x^2 on [0,1]
   n      left     right       mid  trapezoid    Simpson
  10  4.833e-02  5.167e-02  8.333e-04  1.667e-03  1.110e-16
 100  4.983e-03  5.017e-03  8.333e-06  1.667e-05  0.000e+00
1000  4.998e-04  5.002e-04  8.333e-08  1.667e-07  0.000e+00

quad exp(-x^2) = 0.746824132812427   estimated error 8.3e-15
erf reference  = 0.746824132812427   difference +1.1e-16
quad x^2       = 0.33333333333333337

error |Simpson - reference| for exp(-x^2)
  n=  2  3.563e-04
  n=  4  3.125e-05
  n=  8  1.988e-06
  n= 16  1.246e-07
  n= 32  7.795e-09

error |trapezoid - reference| for exp(-x^2)
  n=  10  6.133e-04
  n= 100  6.131e-06
  n=1000  6.131e-08
\end{verbatim}

Four things stand out. The difference halves each time $n$ doubles, consistent with $1/1024$
in the last row. The left and right errors drop by a factor of ten, midpoint and
trapezoid by a factor of a hundred. The Simpson column contains $1.110\times10^{-16}$ and then zero, a residue
of rounding, and Simpson's error for $e^{-x^2}$ drops by about a factor of sixteen
each time $n$ is doubled.

That code does not prove a single theorem. It checks that the formulas of this chapter give
the numbers stated. The proof of integrability of $x^2$ is in the computation $U_n-L_n=1/n$
in the previous section: the output of a program is data, not an argument.

\begin{untukdosen}
The teaching order I suggest: start from the two different answers in the opening, then
the definition of a partition and compute $U$, $L$ for $x^2$ with four slices on the board. After
that the refinement lemma, the criterion, and only then the classes of functions. Do not postpone the fundamental
theorem of calculus: students need to see the hard work of the criterion bear fruit as a computational tool.

Two misconceptions that appear often. First, the integral is taken to be the same as the antiderivative,
so that $F'=f$ is taken to hold by definition, whereas it is a theorem. Second, a function
with many discontinuity points is taken to be certainly not integrable, because the Dirichlet
example is too often used as the only example. Thomae's function removes
that assumption in a single page.

Blackboard activity: ask students to guess $\int_0^1x^2dx$ only from $U_n$ and
$L_n$ for $n=1,2,4$, then to say how many $n$ are needed for the error to be below
$10^{-3}$ without computing the integral. The answer $n>1000$ comes out of $1/n$, and that moment
explains why an antiderivative is more valuable than quadrature. Sign of readiness
to move on: students can explain the difference in direction of the two parts of the fundamental theorem of calculus,
and can name one integrable function that is not continuous.
\end{untukdosen}

\section{Summary}
\tingkat{basic}

\begin{ringkasan}
\begin{itemize}
\item The Riemann integral is the only number trapped between all
      lower sums and all upper sums; when the two brackets meet, that value is called
      the integral.
\item The integrability criterion replaces the examination of infinitely many partitions
      with one inequality $U(\mathcal{P},f)-L(\mathcal{P},f)<\varepsilon$.
\item Continuous functions and monotone functions are integrable; the Dirichlet function is not, because
      every slice of it contains rational and irrational numbers, whereas Thomae's function
      is integrable even though it is broken at every rational number.
\item The accumulation function $F(x)=\int_a^xf$ is always continuous and its derivative is $f$ at every
      point where $f$ is continuous, so every continuous function has an antiderivative.
\item The second part states $\int_a^bF'=F(b)-F(a)$; its proof is the mean value theorem
      plus telescoping, with the condition that $F'$ be integrable.
\item Substitution and integration by parts are the chain rule and the product rule
      read backwards. An improper integral is a limit of Riemann integrals, and the integral
      test links it to the series of Chapter 5.
\end{itemize}
\end{ringkasan}

\section{Exercises}
\tingkat{intermediate}

\begin{latihan}
\begin{enumerate}
  \item Compute $U_n$ and $L_n$ for $f(x)=x$ on $[0,1]$ with a uniform partition,
        express both as fractions in $n$, then determine $U_n-L_n$. Conclude
        that $f$ is integrable and compute the integral.
  \item Prove that $f$ integrable on $[a,b]$ implies $|f|$ integrable
        and $\bigl|\int_a^bf\bigr|\le\int_a^b|f|$. Show that the converse is false with
        one example.
  \item Assess this claim: ``Every bounded function on $[a,b]$ is integrable, since its
        upper sum and lower sum are always finite.'' True, false, or not yet
        enough evidence? Name which part is supported by proof and which part is
        not.
  \item Prove that an increasing function on $[a,b]$ is integrable by showing
        $U(\mathcal{P},f)-L(\mathcal{P},f)\le(b-a)\bigl(f(b)-f(a)\bigr)/n$ for
        a uniform partition into $n$ parts. Explain why that bound does not hold for
        a function that is not monotone.
  \item Compute $\int_0^1(3x^2-2x+1)\,dx$ with the fundamental theorem of calculus, then check
        the result using Simpson's rule with $n=100$ and the function
        \texttt{scipy.integrate.quad}. Report all three numbers and their differences.
  \item Compute $\int_1^{10}x^{-2}dx$ and $\int_1^{\infty}x^{-2}dx$. Report the difference
        between them, then explain why that difference equals $1/10$ without
        recomputing.
  \item Prove that if $f$ is continuous, $f\ge0$ on $[a,b]$, and $\int_a^bf=0$, then
        $f(x)=0$ for every $x\in[a,b]$. Also show that the statement fails
        if $f$ is only integrable.
  \item Compute $\int_0^{\pi/2}\sin x\cos x\,dx$ in two ways: the substitution
        $u=\sin x$, and the identity $\sin2x=2\sin x\cos x$. Compare the results with
        the numerical value from \texttt{quad}.
  \item Show that the Dirichlet function on $[0,1]$ satisfies
        $U(\mathcal{P},D)-L(\mathcal{P},D)=1$ for every partition, then explain
        why that fact is enough to conclude that it is not integrable.
  \item Estimate $\sum_{n=101}^{\infty}1/n^2$ with an integral bound and compare with
        the true value $9.950167\times10^{-3}$; what is the relative error?
\end{enumerate}

\noindent\textbf{Hints.} Item 3: the claim contains two different statements, about
the boundedness of the sums and about the equality of the two integrals. Item 7: use
continuity at one point to make the integral small on a narrow interval around
that point. Item 10: use $\sum_{n=N+1}^\infty1/n^2\le\int_N^\infty x^{-2}dx=1/N$
with $N=100$.
\end{latihan}

\begin{untukpeneliti}
What is genuinely proved in this chapter: the integrability criterion with both of its directions,
the integrability of continuous functions (using uniform continuity from Chapter 11), of monotone
functions, of Thomae's function, both parts of the fundamental theorem of calculus, the mean value theorem
for integrals, and the status of substitution and integration by parts as consequences, not
axioms.

What is quoted without proof and should be called so in class: Lebesgue's criterion
(Rudin, Theorem 11.33(b); Abbott, Chapter 7), Volterra's example of a bounded derivative that
is not integrable, and the absence of an elementary antiderivative for $e^{-x^2}$ (Liouville's
theorem).

Primary sources: Rudin, \emph{Principles of Mathematical Analysis}, Chapter 6, the section
Definition and Existence of the Integral (Theorems 6.4, 6.6, 6.8, 6.9, 6.10) and the section
Integration and Differentiation (6.19--6.22). Rudin sets it up for the
Riemann-Stieltjes integral; the Riemann integral is the case $\alpha(x)=x$. Abbott,
\emph{Understanding Analysis}, Chapter 7, is slower but more readable; Stewart only
for the computational rules.

Open question: how sharp the quadrature error bounds are when only $|f''|$ is
known, and how to estimate the error without any derivative at all; the standard
answers are adaptive quadrature and the difference of two levels of refinement. Limits of validity:
bounded functions on a closed bounded interval, improper integrals as limits, without
measure theory, and vector-valued integrals are not touched.
\end{untukpeneliti}

\peta{Chapter 13 follows Rudin, Chapter 6: the section Definition and Existence of the Integral
(partitions, upper and lower sums, the refinement lemma as Theorem 6.4, the criterion as
6.6, continuous and monotone functions as 6.8 and 6.9) and the section Integration and
Differentiation (6.19 change of variables, 6.21 the fundamental theorem of calculus, 6.22 integration by
parts). In Abbott, Chapter 7, from The Definition of the Riemann Integral to
Lebesgue's Criterion. In Stewart: The Substitution Rule and Average Value of a
Function (Chapter 5), Integration by Parts and Improper Integrals (Chapter 7), and 11.3
The Integral Test, the bridge to Chapter 5 of this book.}

